\SourceLeaf{140}
\Needspace{8\baselineskip}
\begin{center}
\textit{Ambiguous forms.}
\end{center}
\Article{163}
We have already mentioned in passing that a form $F$ may contain another form $F'$ both properly and improperly. This clearly happens if between the forms $F,F'$ we can place another form $G$ so that $F$ contains $G$ and $G$ contains $F'$, and the form $G$ is improperly equivalent to itself. For suppose $F$ contains $G$ properly or improperly. Since $G$ contains $G$ improperly, $F$ also contains $G$ improperly or properly, respectively, and thus in either case both properly and improperly (Article 159). It follows in the same way that however $G$ is assumed to contain $F'$, $F$ must always contain $F'$ both properly and improperly. --- That forms improperly equivalent to themselves exist can be seen from the most obvious case, where the middle term of the form $=0$. Such a form is its own opposite (Article 159) and so is improperly equivalent to itself. More generally, every form $(a,b,c)$ in which $2b$ is divisible by $a$ has this property. For the form $(c,b,a)$ is its neighbor on the left (Article 160), and so is properly equivalent to it. But by Article 159, $(c,b,a)$ is improperly equivalent to the form $(a,b,c)$, so $(a,b,c)$ is improperly equivalent to itself. We will call forms $(a,b,c)$ in which $2b$ is divisible by $a$ \textit{ambiguous}. We therefore have the following theorem:

\textit{A form $F$ contains another form $F'$ both properly and improperly if an ambiguous form can be found that is contained in $F$ and contains $F'$.} But this proposition also has a converse, namely:

\Needspace{18\baselineskip}
\begin{center}
\textit{Theorem on the case where one form is contained in another both properly and improperly.}
\end{center}
\Article{164}
\textsc{Theorem.} \textit{If the form}
\[
Axx+2Bxy+Cyy\quad\mathinner{.\,.\,.\,.\,.}\quad(F)
\]
\textit{contains the form}
\[
A'x'x'+2B'x'y'+C'y'y'\quad\mathinner{.\,.\,.\,.}\quad(F')
\]
\textit{both properly and improperly, an ambiguous form can be found that is contained in $F$ and contains the form $F'$.}

Suppose the form $F$ changes into the form $F'$ both under the substitution
\[
x=\alpha x'+\beta y',\qquad y=\gamma x'+\delta y'
\]
\SourceLeaf{141}
and under the following dissimilar substitution:
\[
x=\alpha'x'+\beta'x',\qquad y=\gamma'x'+\delta'y'
\]
Denoting the numbers $\alpha\delta-\beta\gamma,\ \alpha'\delta'-\beta'\gamma'$ by $e,e'$, we have $B'B'-A'C'=ee(BB-AC)=e'e'(BB-AC)$, and so $ee=e'e'$. Since by hypothesis $e,e'$ have opposite signs, $e=-e'$, or $e+e'=0$. It is now clear that if, in $F'$, we substitute for $x'$ the expression $\delta'x''-\beta'y''$, and make the corresponding replacement $y',-\gamma'x''+\alpha'y''$, the same form results as when we write in $F$
\[
\begin{aligned}
\text{\textit{either} 1) for }x\quad&\alpha(\delta'x''-\beta'y'')+\beta(-\gamma'x''+\alpha'y'')\\
\text{\textit{that is}}\quad&(\alpha\delta'-\beta\gamma')x''+(\beta\alpha'-\alpha\beta')y''\\
\text{and for }y\quad&\gamma(\delta'x''-\beta'y'')+\delta(-\gamma'x''+\alpha'y'')\\
\text{\textit{that is}}\quad&(\gamma\delta'-\delta\gamma')x''+(\delta\alpha'-\gamma\beta')y''\\
\text{\textit{or} 2) for }x\quad&\alpha'(\delta'x''-\beta'y'')+\beta'(-\gamma'x''+\alpha'y'')\quad\text{\textit{that is}}\quad e'x''\\
\text{and for }y\quad&\gamma'(\delta'x''-\beta'y'')+\delta'(-\gamma'x''+\alpha'y'')\quad\text{\textit{that is}}\quad e'y''
\end{aligned}
\]
Thus, denoting the numbers $\alpha\delta'-\beta\gamma',\ \beta\alpha'-\alpha\beta',\ \gamma\delta'-\delta\gamma',\ \delta\alpha'-\gamma\beta'$ by $a,b,c,d$, the form $F$ changes under the two substitutions
\[
x=ax''+by'',\quad y=cx''+dy'';\qquad x=e'x'',\quad y=e'y''
\]
into the same form. This gives the following three equations:
\[
\begin{aligned}
Aaa+2Bac+Ccc&=Ae'e'&\quad\mathinner{.\,.\,.\,.\,.\,.\,.}\quad&[1]\\
Aab+B(ad+bc)+Ccd&=Be'e'&\quad\mathinner{.\,.\,.\,.\,.}\quad&[2]\\
Abb+2Bbd+Cdd&=Ce'e'&\quad\mathinner{.\,.\,.\,.\,.\,.\,.}\quad&[3]
\end{aligned}
\]
From the values of $a,b,c,d$, however, we find
\[
ad-bc=ee'=-ee=-e'e'\quad\mathinner{.\,.\,.\,.\,.\,.}\quad[4]
\]
Then from $d[1]-c[2]$ we get
\[
(Aa+Bc)(ad-bc)=(Ad-Bc)e'e'
\]
and so
\[
A(a+d)=0
\]
Further, from $(a+d)[2]-b[1]-c[3]$ we get

\SourceLeaf{142}

\[
(Ab+B(a+d)+Cc)(ad-bc)=(-Ab+B(a+d)-Cc)e'e'
\]

and so

\[
B(a+d)=0
\]

Finally, from $a[3]-b[2]$ we get

\[
(Bb+Cd)(ad-bc)=(-Bb+Ca)e'e'
\]

and so

\[
C(a+d)=0
\]

Since $A,B,C$ cannot all be $=0$, we must have $a+d=0$, or $a=-d$.

From $a[2]-b[1]$ we get

\[
(Ba+Cc)(ad-bc)=(Ba-Ab)e'e'
\]

and thus

\[
Ab-2Ba-Cc=0\quad\mathinner{.\ .\ .\ .\ .\ .\ .\ .}\quad[5]
\]

The equations $e+e'=0$, $a+d=0$, or

\[
\alpha\delta-\beta\gamma+\alpha'\delta'-\beta'\gamma'=0,\quad\alpha\delta'-\beta\gamma'-\gamma\beta'+\delta\alpha'=0
\]

give $(\alpha+\alpha')(\delta+\delta')=(\beta+\beta')(\gamma+\gamma')$, or

\[
(\alpha+\alpha'):(\gamma+\gamma')=(\beta+\beta'):(\delta+\delta')
\]

Let this ratio\HistoricalNote{*)}{If all of $\alpha+\alpha'$, $\gamma+\gamma'$, $\beta+\beta'$, $\delta+\delta'$ were $=0$, the ratio would be indeterminate, and the method would not apply. But a little thought shows that this is inconsistent with our assumptions. For we would have $\alpha\delta-\beta\gamma=\alpha'\delta'-\beta'\gamma'$, that is, $e=e'$, and thus, since $e=-e'$, $e=e'=0$. Then $B'B'-A'C'$, that is, the determinant of the form $F'$, would also become $=0$. We have excluded such forms entirely.} be written in lowest terms as $m:n$, so that $m,n$ are relatively prime, and choose $\mu,\nu$ so that $\mu m+\nu n=1$. Further, let $r$ be the greatest common divisor of the numbers $a,b,c$. Its square therefore divides $aa+bc$, or $bc-ad$, or $ee$, so $r$ also divides $e$. With these choices, suppose the form $F$ changes under the substitution

\[
x=mt+\frac{\nu e}{r}u,\qquad y=nt-\frac{\mu e}{r}u
\]

into the form $Mtt+2Ntu+Puu$, $(G)$. This form will be ambiguous and will contain the form $F'$.

\SourceLeaf{143}

\emph{Proof.} I. To show that the form $G$ is ambiguous, we will prove that

\[
M(b\mu\mu-2a\mu\nu-c\nu\nu)=2Nr
\]

Since $r$ divides $a,b,c$, the number $\frac{1}{r}(b\mu\mu-2a\mu\nu-c\nu\nu)$ is an integer, and so $2N$ is a multiple of $M$. Now

\[
\begin{gathered}M=Amm+2Bmn+Cnn,\\Nr=(Am\nu-B(m\mu-n\nu)-Cn\mu)e\quad\mathinner{.\ .}\quad[6]\end{gathered}
\]

Further, expansion easily confirms that

\[
\begin{aligned}2e+2a&=e-e'+a-d\\&=(\alpha-\alpha')(\delta+\delta')-(\beta-\beta')(\gamma+\gamma')\end{aligned}
\]

\[
2b=(\alpha+\alpha')(\beta-\beta')-(\alpha-\alpha')(\beta+\beta')
\]

Since $m(\gamma+\gamma')=n(\alpha+\alpha')$, $m(\delta+\delta')=n(\beta+\beta')$, this gives

\[
m(2e+2a)=-2nb\quad\text{or}
\]

\[
me+ma+nb=0\quad\mathinner{.\ .\ .\ .\ .\ .\ .\ .\ .}\quad[7]
\]

Similarly,

\[
\begin{aligned}2e-2a&=e-e'-a+d\\&=(\alpha+\alpha')(\delta-\delta')-(\beta+\beta')(\gamma-\gamma')\end{aligned}
\]

\[
2c=(\gamma-\gamma')(\delta+\delta')-(\gamma+\gamma')(\delta-\delta')
\]

and so $n(2e-2a)=-2mc$, or

\[
ne-na+mc=0\quad\mathinner{.\ .\ .\ .\ .\ .\ .\ .\ .}\quad[8]
\]

Now add to $mm(b\mu\mu-2a\mu\nu-c\nu\nu)$ the expression

\[
\begin{gathered}(1-m\mu-n\nu)(m\nu(e-a)+(m\mu+1)b)\\+(me+ma+nb)(m\mu\nu+\nu)+(ne-na+mc)m\nu\nu\end{gathered}
\]

which clearly $=0$, because

\[
1-\mu m-\nu u=0,\quad me+ma+nb=0,\quad ne-na+mc=0
\]

Expanding the products and removing the terms that cancel gives $2m\nu e+b$. So

\[
mm(b\mu\mu-2a\mu\nu-c\nu\nu)=2m\nu e+b\quad\mathinner{.\ .\ .\ .\ .}\quad[9]
\]

Similarly, adding to $mn(b\mu\mu-2a\mu\nu-c\nu\nu)$ the following:

\SourceLeaf{144}

\[
\begin{gathered}(1-m\mu-n\nu)((n\nu-m\mu)e-(1+m\mu+n\nu)a)\\-(me+ma+nb)m\mu\mu+(ne-na+mc)n\nu\nu\end{gathered}
\]

gives

\[
mn(b\mu\mu-2a\mu\nu-c\nu\nu)=(n\nu-m\mu)e-a\quad .\ .\ .\ .\ [10]
\]

Finally, adding to $nn(b\mu\mu-2a\mu\nu-c\nu\nu)$ the expression

\[
\begin{gathered}(m\mu+n\nu-1)(n\mu(e+a)+(n\nu+1)c)\\-(me+ma+nb)n\mu\mu-(ne-na+mc)(n\mu\nu+\mu)\end{gathered}
\]

gives

\[
nn(b\mu\mu-2a\mu\nu-c\nu\nu)=-2n\mu e-c\quad .\ .\ .\ .\ [11]
\]

Now from $9$, $10$, $11$, we derive

\[
\begin{gathered}(Amm+2Bmn+Cnn)(b\mu\mu-2a\mu\nu-c\nu\nu)\\=2e(Am\nu+B(n\nu-m\mu)-Cn\mu)+Ab-2Ba-Cc\end{gathered}
\]

or, because of $[6],$,

\[
M(b\mu\mu-2a\mu\nu-c\nu\nu)=2Nr.\quad Q.\ E.\ D.
\]

II. To prove that the form $G$ contains the form $F'$, we will show \emph{first} that $G$ changes into $F'$ on putting

\[
\begin{gathered}t=(\mu\alpha+\nu\gamma)x'+(\mu\beta+\nu\delta)y',\\u=\frac re(n\alpha-m\gamma)x'+\frac re(n\beta-m\delta)y'\ .\ .\ .\ (S)\end{gathered}
\]

and \emph{second} that $\frac{r}{e}(n\alpha-m\gamma)$, $\frac{r}{e}(n\beta-m\delta)$ are integers.

1. Since $F$ changes into $G$ on putting

\[
x=mt+\frac{\nu e}{r}u,\quad y=nt-\frac{\mu e}{r}u
\]

the form $G$ changes under substitution $(S)$ into the same form as $F$ does on putting

\[
\begin{aligned}x={}&m((\mu\alpha+\nu\gamma)x'+(\mu\beta+\nu\delta)y')\\&+\nu((n\alpha-m\gamma)x'+(n\beta-m\delta)y')\end{aligned}
\]

\emph{that is}, $=\alpha(m\mu+n\nu)x'+\beta(m\mu+n\nu)y'$, or $=\alpha x'+\beta y'$,

and

\[
\begin{aligned}y={}&n((\mu\alpha+\nu\gamma)x'+(\mu\beta+\nu\delta)y')\\&-\mu((n\alpha-m\gamma)x'+(n\beta-m\delta)y')\end{aligned}
\]

\emph{that is}, $=\gamma(n\nu+m\mu)x'+\delta(n\nu+m\mu)y'$, or $=\gamma x'+\delta y'$.

\SourceLeaf{145}

But this substitution changes $F$ into $F'$, so substitution $(S)$ also changes $G$ into $F'$.

2. From the values of $e,b,d$ we find $\alpha'e+\gamma b-\alpha d=0$, or, because $d=-a$, $n\alpha'e+n\alpha a+n\gamma b=0$. Then $[7]$ gives $n\alpha'e+n\alpha a=m\gamma e+m\gamma a$, or

\[
(n\alpha-m\gamma)a=(m\gamma-n\alpha')e\quad .\ .\ .\ .\ .\ .\ .\ [12]
\]

Further, $\alpha nb=-\alpha m(e+a)$, $\gamma mb=-m(\alpha'e+\alpha a)$, and so

\[
(n\alpha-m\gamma)b=(\alpha'-\alpha)me\quad .\ .\ .\ .\ .\ .\ .\ [13]
\]

Finally, $\gamma'e-\gamma a+\alpha c=0$. Multiplying by $n$ and substituting for $na$ its value from $[8]$ gives

\[
(n\alpha-m\gamma)c=(\gamma-\gamma')ne\quad .\ .\ .\ .\ .\ .\ .\ [14]
\]

Similarly, we derive $\beta'e+\delta b-\beta d=0$, or $n\beta'e+n\delta b+n\beta a=0$, and therefore, by $[7],$, $n\beta'e+n\beta a=m\delta e+m\delta a$, or

\[
(n\beta-m\delta)a=(m\delta-n\beta')e\quad .\ .\ .\ .\ .\ .\ .\ [15]
\]

Further, $\beta nb=-\beta m(e+a)$, $\delta mb=-m(\beta'e+\beta a)$, and so

\[
(n\beta-m\delta)b=(\beta'-\beta)me\quad .\ .\ .\ .\ .\ .\ .\ [16]
\]

Finally, $\delta e-\delta a+\beta c=0$. Multiplying by $n$ and substituting for $na$ its value from $[8]$ gives

\[
(n\beta-m\delta)c=(\delta-\delta')ne\quad .\ .\ .\ .\ .\ .\ .\ [17]
\]

Now, since the greatest common divisor of the numbers $a,b,c$ is $r$, integers $\mathfrak{A},\mathfrak{B},\mathfrak{C}$ can be chosen so that

\[
\mathfrak{A}a+\mathfrak{B}b+\mathfrak{C}c=r
\]

Then $12$, $13$, $14$ and $15$, $16$, $17$ give

\[
\begin{aligned}&\mathfrak A(m\gamma-n\alpha')+\mathfrak B(\alpha'-\alpha)m+\mathfrak C(\gamma-\gamma')n\\&\qquad=\frac re(n\alpha-m\gamma)\\&\mathfrak A(m\delta-n\beta')+\mathfrak B(\beta'-\beta)m+\mathfrak C(\delta-\delta')n\\&\qquad=\frac re(n\beta-m\delta)\end{aligned}
\]

so $\frac{r}{e}(n\alpha-n\gamma),\quad\frac{r}{e}(n\beta-m\delta)$ are integers. $Q.\ E.\ D.$

\SourceLeaf{146}

\Article{165}

\emph{Example.} The form $3xx+14xy-4yy$ changes into the form $-12x'x'-18x'y'+39y'y'$ both properly, on putting

\[
x=4x'+11y',\quad y=-x'-2y'
\]

and improperly, on putting

\[
x=-74x'+89y',\quad y=15x'-18y'
\]

Here $\alpha+\alpha'$, $\beta+\beta'$, $\gamma+\gamma'$, $\delta+\delta'$ are $-70$, $100$, $14$, $-20$, and $-70:14=100:-20=5:-1$. So we put $m=5$, $n=-1$, $\mu=0$, $\nu=-1$. The numbers $a,b,c$ are found to be $-237$, $-1170$, $48$, whose greatest common divisor $=3=r$; finally, $e=3$. Thus transformation $(S)$ is $x=5t-u$, $y=-t$. It changes the form $(3,7,-4)$ into the ambiguous form $tt-16tu+3uu$.

If the forms $F,F'$ are equivalent, the form $G$, contained in $F$, is also contained in $F'$. But since it also contains that same form, it is equivalent to it, and so also to the form $F$. In this case, therefore, the theorem can be stated as follows:

\emph{If $F,F'$ are both properly and improperly equivalent, an ambiguous form equivalent to both can be found.} --- In this case $e=\pm1$, so $r.$, which divides $e$, must also be $=1$.

Let this suffice for the transformation of forms in general. We now turn to \emph{representations}.
